% NOETHER PAPER 4 — KOREAN TRANSLATION-PRODUCER DRAFT — UNCHECKED
% Unit: T04-U20; current-authority whole lines 3961--3973
% Source slice: 1,586 LF UTF-8 bytes; SHA-256 B5C21677916BEFF38DE0ED3FD59FEBE6E5DA413ACC8B255BFA6F524113188BFB
% Pointer: NOETH-DE-AUTH-v004-20260804; SHA-256 A1C62FDACAA34DFC1B806DC18258F2E732539F7AE9D85AA4BD9E1067B8749D9F
% This is producer translation only: no source/Korean/formula review, compilation, or rendering.

경우 2)에서도 연산 (23.)은 $\tau$번 적용할 수 있지만, 연산 (24.)는 $\rho$번째 뒤에 중단된다. 따라서 (25.)의 종결식에서는 첫째 합 기호를 $\sum_{\alpha=0}^{\rho}$로 바꾸어야 한다. 경우 3)과 4)에서는 연산 (23.)을 $\sigma$번만 실행할 수 있다. 이때 (24.) 대신 다음 관계가 들어간다.
\begin{equation}
\begin{aligned}
\pair{q_\rho A^\sigma}{p_\tau B_{\rho+\sigma-\tau}}
&=(-1)^{\rho\sigma}
 \pair{q_{\rho-(\tau-\sigma)}^{(\sigma+1\ldots\tau)}B^{n-\rho+\tau-\sigma}}{A_{n-\sigma}p_{\tau-(\tau-\sigma)}^{(\sigma+1\ldots\tau)}}\\
&\quad+\sum_{i_1=1}^{\sigma}(-1)^{\rho(i_1-1)}
 \pair{q_{\rho-1}^{(i_1)}[A_{n-\sigma}y^{(1)}\cdots y^{(i_1-1)}]^{\sigma-i_1+1}}
 {y^{(i_1+1)}\cdots y^{(\tau)}B_{\rho+\sigma-\tau}};
\end{aligned}
\tag{28}
\end{equation}
(28.)을 $(\tau-\sigma)$번 반복하면, 그동안 합 기호는 $\sum_{i_\beta=i_{\beta-1}+1}^{\sigma+\beta-1}$로 바뀌며, (25.)에서 지수 $\alpha=\tau-\sigma$에 해당하는 개별 합의 모든 항을 각각 알맞은 부호와 함께 얻는다. 이는 (28.)의 합 기호 아래 항들에 (23.)을 $(\sigma-i_1+1)$번 적용하면 알 수 있다. 그 뒤에는 경우 3)과 4)가 경우 1)과 2)로 넘어간다. 실제로 $\sigma,\tau$에 대응하는 수들은 $\sigma'=\sigma-i_{\tau-\sigma}+\tau-\sigma$, $\tau'=\tau-i_{\tau-\sigma}=\sigma'$가 되며, 절차를 계속하면 $\tau'<\sigma'$가 된다. 그러므로 (25.)의 종결식에서는 첫째 합 기호를 $\sum_{\alpha=\tau-\sigma}^{\tau,\rho}$로 바꾸어야 한다.
